% from ega1-10-fr.tex
\oldpage[I]{210}
est le préschéma formel limite inductive des $Y_n$, et comme les conditions de
(\hyperref[I.10.6.4-fr]{10.6.4}) sont satisfaites, cela prouve que ce préschéma
formel est localement noethérien, et noethérien si $\mathfrak{X}$ l'est (puisque
alors $Y_0$ l'est en vertu de (\hyperref[I.6.1.4-fr]{6.1.4})).

\begin{definition}[10.14.2]
\label{I.10.14.2-fr}
On appelle \emph{sous-préschéma fermé} d'un préschéma formel $\mathfrak{X}$ tout
préschéma formel
$(\mathfrak{Y},(\mathscr{O}_{\mathfrak{X}}/\mathscr{A})|\mathfrak{Y})$ où
$\mathscr{A}$ est un $\mathscr{O}_{\mathfrak{X}}$-Module cohérent~; on dit que ce
préschéma est le sous-préschéma fermé défini par $\mathscr{A}$.
\end{definition}

Il est clair que la correspondance ainsi définie entre
$\mathscr{O}_{\mathfrak{X}}$-Modules cohérents et sous-préschémas fermés de
$\mathfrak{X}$ est biunivoque.

Le morphisme d'espaces topologiquement annelés
$j=(\psi,\theta):\mathfrak{Y}\to\mathfrak{X}$, où $\psi$ est l'injection
$\mathfrak{Y}\to\mathfrak{X}$ et $\theta^\sharp$ l'homomorphisme canonique
$\mathscr{O}_{\mathfrak{X}}\to\mathscr{O}_{\mathfrak{X}}/\mathscr{A}$, est
évidemment (\hyperref[I.10.4.5-fr]{10.4.5}) un morphisme de préschémas formels,
qu'on appelle l'\emph{injection canonique} de $\mathfrak{Y}$ dans $\mathfrak{X}$.
On notera que si $\mathfrak{X}=\operatorname{Spf}(A)$, où $A$ est adique
noethérien, on a $\mathscr{A}=\mathfrak{a}^\Delta$, où $\mathfrak{a}$ est un
idéal de $A$ (\hyperref[I.10.10.5-fr]{10.10.5}), et il résulte aussitôt de ce
qui précède que l'on a alors
$\mathfrak{Y}=\operatorname{Spf}(A/\mathfrak{a})$ à un isomorphisme près, et que
$j$ correspond (\hyperref[I.10.2.2-fr]{10.2.2}) à l'homomorphisme canonique
$A\to A/\mathfrak{a}$.

On dit qu'un morphisme $f:\mathfrak{Z}\to\mathfrak{X}$ de préschémas formels
localement noethériens est une \emph{immersion fermée} si elle se factorise en
$\mathfrak{Z}\xrightarrow{g}\mathfrak{Y}\xrightarrow{j}\mathfrak{X}$, où $g$
est un isomorphisme de $\mathfrak{Z}$ sur un sous-préschéma fermé
$\mathfrak{Y}$ de $\mathfrak{X}$ et $j$ l'injection canonique. Comme $j$ est un
monomorphisme d'espaces annelés, $g$ et $\mathfrak{Y}$ sont nécessairement
\emph{uniques}.

\begin{proposition}[10.14.3]
\label{I.10.14.3-fr}
Une immersion fermée est un morphisme de type fini.
\end{proposition}

On se ramène aussitôt au cas où $\mathfrak{X}$ est un schéma formel affine
$\operatorname{Spf}(A)$ et $\mathfrak{Y}=\operatorname{Spf}(A/\mathfrak{a})$~;
la proposition résulte de (\hyperref[I.10.13.1-fr]{10.13.1, c)}).

\begin{lemma}[10.14.4]
\label{I.10.14.4-fr}
Soit $f:\mathfrak{Y}\to\mathfrak{X}$ un morphisme de préschémas formels
localement noethériens, et soit $(U_\alpha)$ un recouvrement de
$f(\mathfrak{Y})$ par des ouverts formels affines noethériens de
$\mathfrak{X}$, tels que les $f^{-1}(U_\alpha)$ soient des ouverts formels
affines noethériens de $\mathfrak{Y}$. Pour que $f$ soit une immersion fermée,
il faut et il suffit que $f(\mathfrak{Y})$ soit une partie fermée de
$\mathfrak{X}$ et que, pour tout $\alpha$, la restriction de $f$ à
$f^{-1}(U_\alpha)$ corresponde (\hyperref[I.10.4.6-fr]{10.4.6}) à un
homomorphisme surjectif
$\Gamma(U_\alpha,\mathscr{O}_{\mathfrak{X}})\to
\Gamma(f^{-1}(U_\alpha),\mathscr{O}_{\mathfrak{Y}})$.
\end{lemma}

Les conditions sont évidemment nécessaires. Inversement, si elles sont
remplies, et si on désigne par $\mathfrak{a}_\alpha$ le noyau de
$\Gamma(U_\alpha,\mathscr{O}_{\mathfrak{X}})\to
\Gamma(f^{-1}(U_\alpha),\mathscr{O}_{\mathfrak{Y}})$, on définit un faisceau
cohérent d'idéaux $\mathscr{A}$ de $\mathscr{O}_{\mathfrak{X}}$ en prenant
$\mathscr{A}|U_\alpha=\mathfrak{a}_\alpha^\Delta$, et en prenant $\mathscr{A}$
nul dans le complémentaire de la réunion des $U_\alpha$. En effet, puisque
$f(\mathfrak{Y})$ est fermé et que le support de
$\mathfrak{a}_\alpha^\Delta$ est $U_\alpha\cap f(\mathfrak{Y})$, tout revient
à vérifier que $\mathfrak{a}_\alpha^\Delta$ et
$\mathfrak{a}_\beta^\Delta$ induisent le même faisceau sur un ouvert formel
affine noethérien $V\subset U_\alpha\cap U_\beta$. Or, la restriction de $f$
à $f^{-1}(U_\alpha)$ étant une immersion fermée de ce préschéma formel dans
$U_\alpha$, $f^{-1}(V)$ est un ouvert formel affine noethérien dans
$f^{-1}(U_\alpha)$ et la restriction de $f$ à $f^{-1}(V)$ est une immersion
fermée~; si $\mathfrak{b}$ est le noyau de l'homomorphisme surjectif
$\Gamma(V,\mathscr{O}_{\mathfrak{X}})\to
\Gamma(f^{-1}(V),\mathscr{O}_{\mathfrak{Y}})$ correspondant à cette
restriction, il est immédiat (\hyperref[I.10.10.2-fr]{10.10.2}) que
$\mathfrak{a}_\alpha^\Delta$ induit $\mathfrak{b}^\Delta$ sur $V$. Le faisceau
d'idéaux $\mathscr{A}$ étant ainsi défini, il est alors clair que
$f=g\circ j$, où $j:\mathfrak{Z}\to\mathfrak{X}$ est l'injection canonique du
sous-préschéma fermé $\mathfrak{Z}$ de $\mathfrak{X}$ défini par
$\mathscr{A}$, et $g$ un isomorphisme de $\mathfrak{Y}$ sur $\mathfrak{Z}$.

\begin{proposition}[10.14.5]
\label{I.10.14.5-fr}
\begin{enumerate}
  \item[\rm{(i)}] Si $f:\mathfrak{Z}\to\mathfrak{Y}$,
  $g:\mathfrak{Y}\to\mathfrak{X}$ sont des immersions fermées de préschémas
  formels localement noethériens, $g\circ f$ est une immersion fermée.

